\documentclass{standalone}
\usepackage{circuitikz}
\usetikzlibrary{calc}
\definecolor{circuitnotemark}{HTML}{FE00FE}
\makeatletter
\pgfdeclareshape{smacoax}{
  \anchor{center}{\pgfpointorigin}
  \anchor{text}{\pgfpointorigin}
  \anchor{north}{\pgfpoint{0cm}{0.3cm}}
  \anchor{south}{\pgfpoint{0cm}{-0.3cm}}
  \anchor{east}{\pgfpoint{0.3cm}{0cm}}
  \anchor{west}{\pgfpoint{-0.3cm}{0cm}}
  \anchor{pin 1}{\pgfpoint{-0.7cm}{0cm}}
  \anchor{pin 2}{\pgfpoint{0cm}{-0.7cm}}
  \anchor{bpin 1}{\pgfpoint{-0.3cm}{0cm}}
  \anchor{bpin 2}{\pgfpoint{0cm}{-0.3cm}}
  \backgroundpath{
    \pgfpathmoveto{\pgfpoint{-0.274cm}{-0.122cm}}
    \pgfpatharc{204}{516}{0.3cm}
    \pgfpathmoveto{\pgfpoint{-0.7cm}{0cm}}\pgfpathlineto{\pgfpointorigin}
    \pgfpathmoveto{\pgfpoint{0cm}{-0.7cm}}\pgfpathlineto{\pgfpoint{0cm}{-0.3cm}}
  }
  \foregroundpath{
    \pgfpathcircle{\pgfpointorigin}{0.07cm}
    \pgfusepath{fill}
  }
}
\makeatother
\begin{document}
\begin{circuitikz}[american, line width=0.8pt]
\ctikzset{bipoles/length=1.2cm}
\ctikzset{grounds/scale=1.36}
\coordinate (x2y2) at (2,-2);
\coordinate (x3y2) at (4,-2);
\coordinate (x5y2) at (8,-2);
\coordinate (x7y2) at (12,-2);
\coordinate (x9y2) at (16,-2);
\coordinate (x2y3) at (2,-4);
\coordinate (x9y3) at (16,-4);
\node[smacoax, draw, xscale=-1] (part-J1) at (x2y2) {}; % line 3
\node[anchor=south] at (part-J1.north) {$J_{1}$}; % line 3
\draw (x3y2) to[R, l_=$R_{1}$, a^=$40\,\Omega$] (x5y2); % line 4
\draw (x5y2) to[L, l_=$L_{1}$, a^=$250\,\mathrm{n}\mathrm{H}$] (x7y2); % line 5
\draw (x7y2) to[C, l_=$C_{1}$, a^=$5\mbox{.}6\,\mathrm{p}\mathrm{F}$] (x9y2); % line 6
\node[ground] at (x2y3) {}; % line 7
\node[ground] at (x9y3) {}; % line 8
\draw (part-J1.pin 1) -- (x3y2); % line 10
\draw (x9y2) -- (x9y3); % line 11
\draw (part-J1.pin 2) -- (x2y3); % line 12
\path (1.619,-0.593) rectangle (2.381,0.593); % line 14
\node[anchor=west, inner sep=0, circuitnotemark, font=\large] at (2,0) {X}; % line 14
\node[anchor=south west, inner sep=2pt, yshift=4pt, circuitnotemark, font=\LARGE\bfseries] (circuittitle) at (current bounding box.north west) {X};
\path (circuittitle.south west) rectangle ++(11.034,0.853);
\node[anchor=north east, inner sep=2pt, circuitnotemark, font=\large] (circuitstamp) at ([yshift=-0.593cm]current bounding box.south east) {X};
\path (circuitstamp.north east) rectangle ++(-5.08,-0.593);
\end{circuitikz}
\end{document}
